\section{Discussion and limitations}\label{sec:discussion}
\paragraph{Benchmark}
The four methods differ in one design decision: whether the observation enters the estimate through a likelihood and importance weights, or through a learned update that is trained end to end on complete filtering rollouts. The benchmark separates the two regimes cleanly. When the dynamics are simple enough to learn from clean state pairs, weighting achieves higher accuracy: on the pendulum and Lorenz--96 the best learned particle filter is 9\% and 22\% more accurate than \method. The model-based filters show that the remaining gap to perfect inference is small on the pendulum and dominated by chaotic error amplification on Lorenz--96. When the dynamics are complex (such ascontact-rich) and not governed by a closed-form model, as on the two robot datasets, the learned update excels by a wide margin (25\% and 38\% lower RMSE on ALOHA and REASSEMBLE), with calibrated intervals and about half the latency. The reason is that the filtering distribution of \method is trained on the scored quantity over the whole window, which reduces the modeling errors of the transition and observation models.

\paragraph{The likelihood is the weak point of learned particle filters.}
The noise-recovery analysis (Section~\ref{sec:experiments-noise}) shows all three learned measurement densities converge to the moment-matched Gaussian of the true mixture noise. Maximum-likelihood training of conditional normalizing flows causes the underfitting of the heavy tails, and only the variance is matched instead. \method applies additional energy-score objective contributes to a precise fit of the implied residuals. 

\paragraph{Coverage.}
Coverage probability of an amortized filter relies entirely on the training objective. Two terms are essential: the learned per-coordinate noise scale enabling arbitrary noise scales, and an energy-score based rollout term ensuring proper calibration. Where this balance holds, the intervals are calibrated at every nominal level (Figure~\ref{fig:calibration}). 

\paragraph{Limitations.}
\method is under-covered on the pendulum (0.83 at the 90\% level), mainly in the angle coordinate during the first steps after the known initial state. The learned noise scale and the energy-score term calibrate the other three datasets but do not guarantee calibration. Partially observed chaotic dynamics remain difficult: on Lorenz--96 the six unobserved coordinates can be inferred only through the dynamics, and without a transition model and importance weights \method reaches 0.29 against 0.019 for the reference particle filter. Training also requires careful tuning. The objective of Equation~\eqref{eq:normflowsde_loss} balances four terms, and coverage is sensitive to the rollout weight $\lambda_{\mathrm{rollout}}$, the state scale $\sigma_{\mathrm{state}}$, and the initial noise scale. These were selected on validation trajectories separately for each dataset, and a new system requires the same search.

\section{Conclusion}
\method replaces the likelihood and importance weights of a particle filter by an observation-conditioned normalizing-flow update that is trained end to end on complete filtering rollouts. The benchmark of four datasets with heavy-tailed, multimodal measurement noise shows what this design buys and what it costs. On the two robot datasets it is the most accurate filter with calibrated intervals that are also the narrowest. Because one sample costs one network evaluation per step, without per-particle likelihoods or resampling, it is the fastest method on every dataset, at 1.5 to 2 times the speed of the fastest particle filter. It is also the only learned method whose implied residuals recover the mixture noise, whereas the learned likelihoods of all three particle filters collapse to a Gaussian. The price is paid on simple simulated dynamics, where a learned transition model combined with Bayesian weighting is more accurate, most clearly under partial observation of the chaotic Lorenz--96 system. These results argue for amortized filtering wherever the dynamics are hard to model, the noise is not Gaussian, and inference must be fast and calibrated, and they identify the learned measurement likelihood as the component of particle filters most in need of a better training objective.